\documentclass[11pt,a4paper]{article}

\usepackage[margin=0.8in]{geometry}
\usepackage{setspace}
\usepackage{fancyhdr}
\usepackage{xcolor}
\usepackage{titlesec}

\definecolor{darkblue}{RGB}{35,65,105}

\setstretch{1.05}

\pagestyle{fancy}
\fancyhf{}
\lhead{\textbf{Yashjot Kaur}}
\rhead{\textcolor{darkblue}{Week 2 Journal}}
\cfoot{\thepage}

\renewcommand{\headrulewidth}{0.5pt}

\titleformat{\section}
{\large\bfseries\color{darkblue}}
{}{0em}{}

\begin{document}

\begin{center}
    {\LARGE \textbf{\textcolor{darkblue}{WEEK 2 JOURNAL}}}\\[5pt]
    {\large \textbf{AI-based Real-Time Sensitive Data Detection and Masking}}\\[6pt]
    \textbf{Name: Yashjot Kaur \quad | \quad Roll No.: 1024170452}
\end{center}

\vspace{4pt}
\hrule
\vspace{8pt}

\section*{1. Objective}

The main objective of Week 2 was to understand the working of the project
in more detail and decide how sensitive information would be identified
and protected.

\section*{2. Work Done}

During this week, I explored different examples of sensitive information
such as email IDs, phone numbers, names, addresses, and other personal
details. I studied how these details can appear in normal text and how the
system should recognise them.

I also explored different ways of detecting sensitive information and
understood the difference between information that can be identified using
fixed patterns and information that requires understanding of the text.

\section*{3. My Contribution}

I worked on preparing sample inputs for the project and discussed how
different types of sensitive information should be handled. I also helped
in organising the detection process into simple steps so that it could be
implemented easily in the next stage.

Along with the technical work, I started organising the project
documentation and collected important points to be included in the final
report.

\section*{4. Learning Outcomes}

This week helped me understand the project from an implementation point
of view. I learned that sensitive information can occur in many different
forms and that the system needs to handle these variations carefully.

I also understood the importance of testing the system with realistic
examples before moving towards the final implementation.

\section*{5. Plan for Week 3}

In Week 3, I planned to start the actual detection implementation, test it
with sample data, and work on the initial masking process while continuing
the project report.

\vfill

\begin{center}
    \textcolor{darkblue}{\rule{0.35\textwidth}{0.5pt}}\\[3pt]
    \textit{\small End of Week 2}
\end{center}

\end{document}